\lecture{11}{Wed 19 Feb 2025 13:25}{Ideals and Isomorphism Theorems}
\begin{proposition}\label{prp:2.5}
	Let $I,J$ be ideals of $R$. Then $I+J$ and $I \cap J$ are ideals of $R$.
\end{proposition}
\begin{proof}
	ive alr done this twice wtf idc
\end{proof}
\begin{theorem}[Second Isomorphism Theorem]\label{thm:2.2}
	Let $I,J$ be ideals of a ring $R$. Then
	\[
		\frac{I+J}{J} \cong \frac{I}{I \cap J} 
	.\]
\end{theorem}
\begin{proof}
	owo
\end{proof}
\begin{theorem}[Third Isomorphism Theorem]\label{thm:2.3}
	Let $R$ be a ring with ideals $I,J$ such that $I\subseteq J$. Then $J/I$ is an ideal of $R/I$, and
	\[
		\frac{R/I}{J/I} \cong \frac{R}{J} 
	.\]
\end{theorem}
\begin{proof}
	Follows easily from the first isomorphism theorem.
\end{proof}
\section{Prime and Maximal Ideals}
We spent all of class motivating it.
\begin{definition}[Prime/Maximal Ideals]\label{dfn:2.5}
	Let $R$ be a ring, and $I$ a proper ideal of $R$. Then $I$ is prime if and only if
	\[
		ab\in I \implies a\in I \lor b\in I
	.\]
	A maximal ideal is an ideal $I$ such that if there exists an ideal $J$ with $I\subseteq J$, then $J=R$.
\end{definition}

